\chapter{学习地图：把程序、机器、系统和性能连成一条主线}

\section{本章不是普通导读}

如果一本系统性能书只在开头列出“C++、汇编、CPU、操作系统和性能工具”，它并没有真正降低学习难度。初学者最缺的通常不是名词清单，而是一张可以反复返回的结构图：当你在本册看到 \instruction{mov}、栈、ABI 和 benchmark，在第二册看到 cache line、PMU event、SIMD 或 softmax 时，必须知道这些内容如何共同描述同一个问题：程序怎样在真实机器上运行，并怎样被可信地测量。

本章的目标就是建立这张地图。

刚学完 C/C++ 的读者通常已经知道变量、函数、数组、指针、结构体、类和简单的标准库调用。但从性能工程角度看，这些只是最上层的表达。程序真正运行时，下面还存在许多层：

\begin{center}
\begin{minipage}{0.9\linewidth}
\begin{verbatim}
C/C++ source
  -> language semantics
  -> compiler frontend
  -> IR and optimization
  -> target-specific backend
  -> assembly
  -> machine-code bytes
  -> object file and executable
  -> loader and process
  -> virtual memory
  -> ISA architectural state
  -> microarchitectural execution
  -> cache / TLB / branch predictor / memory system
  -> measured time and performance counters
\end{verbatim}
\end{minipage}
\end{center}

性能优化的难点在于：每一层抽象都在隐藏下一层的复杂性，但这种隐藏从来不是绝对的。日常业务开发依赖抽象来提高效率；系统分析、性能测量和高性能 kernel 开发则经常需要穿过抽象，检查它在什么条件下可靠、在什么边界上泄漏，以及编译器、CPU 或操作系统是否以另一种形式重新解释了你的代码。

\begin{keyidea}
本书的主线不是先堆概念再谈优化，而是从一个 C/C++ 表达式出发，追踪它如何变成机器码、如何改变寄存器和内存、如何穿过 CPU 流水线和缓存层级、如何受到操作系统影响，最后如何形成可信或不可信的性能数据。
\end{keyidea}

\section{一个贯穿全书的例子}

先看一个最小但足够典型的函数：

\begin{lstlisting}[language=C++]
int sum(int const* a, int n)
{
    int s = 0;
    for (int i = 0; i < n; ++i) {
        s += a[i];
    }
    return s;
}
\end{lstlisting}

从 C++ 层看，它的语义直接：遍历数组，把元素累加起来。但要解释它在真实机器上的性能，至少需要回答下面这些问题：

\begin{itemize}
  \item \code{int const* a} 是什么？它只是一个地址，还是还携带了类型和别名语义？
  \item \code{a[i]} 如何变成地址计算？为什么是 \code{base + i * 4}？
  \item \code{s += a[i]} 在汇编中会使用哪些寄存器？会不会每次都访问内存？
  \item 编译器能否把循环展开？能否向量化？为什么有时候不能？
  \item 如果 \code{n} 很小，循环开销和函数调用开销是否比内存访问还重要？
  \item 如果 \code{n} 很大，性能瓶颈是整数加法，还是内存带宽？
  \item 如果 \code{a} 没有按 cache line 对齐，是否一定慢？
  \item 如果多个线程同时处理数组，是否会发生 false sharing？
  \item benchmark 的结果为什么每次都不一样？操作系统调度、CPU 频率、cache 状态分别可能贡献什么噪声？
\end{itemize}

这些问题都不能只靠 C++ 语法解决。它们分别落在语言语义、编译器、汇编、ISA、微架构、操作系统和测量方法上；任何一层被忽略，性能解释都会失去关键依据。

\begin{mentalmodel}
以后看到任何性能代码，都先把它拆成四个视角：

\begin{enumerate}
  \item 源码视角：程序员写了什么语义。
  \item 编译器视角：哪些语义限制了优化，哪些语义允许优化。
  \item 机器视角：实际执行了哪些指令，访问了哪些地址。
  \item 测量视角：数据是否真的证明了你的判断。
\end{enumerate}
\end{mentalmodel}

\section{算法复杂度不是机器成本}

很多读者已经学过算法复杂度，会用 $O(n)$、$O(n \log n)$、$O(1)$ 描述算法规模变化。这当然重要，但它不是性能工程的全部。复杂度告诉我们输入规模变大时增长趋势如何，却不直接告诉我们某个真实程序在某台机器上花多少时间，也不告诉我们瓶颈在哪里。

例如两个函数都是线性扫描数组：

\begin{lstlisting}[language=C++]
for (std::size_t i = 0; i < n; ++i) {
    s += a[i];
}

for (std::size_t i = 0; i < n; ++i) {
    s += a[index[i]];
}
\end{lstlisting}

它们从复杂度看都是 $O(n)$，但机器成本可能完全不同。第一个循环连续读取，容易利用 cache line 和硬件预取；第二个循环间接访问，地址由另一个数组决定，可能造成大量随机 load。两者的比较不能只停在算法层，必须进入地址模式、cache、TLB、乱序执行和带宽模型。

再看两个 $O(1)$ 操作：整数加法和一次系统调用都可以说是常数时间，但它们的常数相差巨大。一个普通整数加法可能只占用很少的核心资源，而系统调用需要从用户态进入内核态，可能涉及权限切换、内核路径、调度和设备状态。复杂度把它们都压缩成常数，却不能替代机器成本分析。

因此，本书使用两套问题：

\begin{center}
\begin{tabular}{p{0.28\linewidth}p{0.55\linewidth}}
\toprule
算法问题 & 机器问题 \\
\midrule
复杂度是什么 & 每次迭代执行哪些指令 \\
输入规模如何增长 & 访问的是连续内存还是随机地址 \\
数据结构抽象是否合适 & 对象布局、padding 和 cache line 如何影响访问 \\
操作次数大概多少 & 指令依赖、分支预测和 load/store 是否限制吞吐 \\
理论上是否最优 & 编译器是否生成预期机器码，测量是否可信 \\
\bottomrule
\end{tabular}
\end{center}

这不是说复杂度不重要，而是说复杂度只是第一层筛选。性能工程要在复杂度可接受的方案里继续比较常数、访存模式、编译器可优化性、并行性和测量稳定性。很多工业级优化并不是把 $O(n)$ 改成 $O(\log n)$，而是让同样 $O(n)$ 的循环更贴近硬件。

\begin{keyidea}
算法复杂度回答“规模增长会怎样”，机器成本回答“这台机器实际怎样执行”。高质量性能分析必须同时包含两者。
\end{keyidea}

\section{抽象层错位：许多误判从这里开始}

学习底层时，最常见的错误不是完全不知道某个术语，而是在一个层次提出问题，却用另一个层次的直觉回答。

例如，问“这个 C++ 变量在哪里”是源码和调试信息层的问题。优化后，这个变量可能在寄存器里，可能被拆成多个值，可能被常量传播消失，甚至根本没有一个稳定位置。如果你用 \code{-O0} 的栈上局部变量经验去解释 \code{-O3} 热路径，就会误判。

再比如，问“这条 load 为什么慢”是微架构和数据状态问题。只回答“因为它是一条内存访问指令”不够。它可能 L1 命中，也可能 L3 miss，也可能被前面的地址依赖卡住，也可能被 TLB miss 或 page fault 影响。指令助记符相同，机器状态不同，成本就不同。

再比如，问“这个 benchmark 为什么快”是测量问题。只回答“因为算法更好”也不够。它可能真的更好，也可能是编译器删除了循环，可能是输入太小全部在 cache 中，可能是只测了 warm cache 场景，可能是没有校验结果导致错误输出。

本书要求你在每次分析时先标注问题层次：

\begin{itemize}
  \item 语言层：程序是否有定义，表达式语义是什么。
  \item 编译器层：优化器能否重写，是否内联、展开、向量化。
  \item ABI 层：参数、返回值、栈和寄存器责任如何约定。
  \item ISA 层：指令如何改变架构状态。
  \item 微架构层：依赖、cache、预测和资源如何影响速度。
  \item 操作系统层：进程、虚拟内存、调度和权限如何参与。
  \item 测量层：数据是否能支撑结论。
\end{itemize}

只要层次标清楚，很多争论会自动变简单。比如“这个函数有没有调用”要看最终二进制，不看源码感觉；“这个地址是否合法”要看进程虚拟地址空间和对象生命周期；“这个优化是否正确”要先看 C++ 语义，而不是只看当前机器上似乎能跑。

\section{计算机系统为什么要分层}

现代计算机太复杂，不可能让程序员直接控制每个晶体管、每条内部总线、每个缓存替换决策和每个乱序执行细节。因此系统必须分层。分层的作用是让上层只依赖下层的一部分承诺。

比如 C++ 不要求你知道整数加法电路如何实现，但它要求机器能表示某些整数值。汇编不要求你知道 CPU 内部有多少执行端口，但它要求 CPU 能按照 ISA 的语义执行 \instruction{add}。操作系统不要求应用知道物理内存条的真实地址，但它提供虚拟地址空间，让进程以为自己拥有连续内存。

分层带来的好处是可编程性，代价是性能信息被藏起来。性能工程正是在这些隐藏处工作。

\subsection{物理层：信号、晶体管和时钟}

最底层是物理世界。电路通过电压高低表示 0 和 1。晶体管可以被看作受控制的开关。大量晶体管组合成逻辑门、触发器、寄存器、加法器、选择器、缓存 SRAM 单元和控制逻辑。

本书不会把你训练成芯片设计工程师，但你必须知道两个事实。

第一，硬件不是“瞬间完成”的。信号传播需要时间，门电路切换需要时间，存储单元读写需要时间。因此 CPU 有时钟周期，有关键路径，有流水线。性能中的 latency 和 throughput，最终都与“工作需要经过多少硬件阶段”和“硬件能否并行处理多个工作”有关。

第二，硬件资源是有限的。一个核心不可能无限发射指令，不可能无限读取内存，不可能无限预测分支，不可能无限保存乱序执行中的临时状态。因此现代 CPU 性能常常不是由某一条高级语句决定，而是由执行端口、load/store 单元、ROB 容量、物理寄存器、L1/L2/L3 cache、TLB、内存带宽等资源共同决定。

\subsection{数字逻辑层：组合逻辑和时序逻辑}

组合逻辑的输出只由当前输入决定，比如加法器、比较器、多路选择器。时序逻辑会记住状态，比如寄存器、触发器、计数器和缓存 tag。CPU 可以被看作一个巨大的状态机：每个时钟周期，它根据当前状态和输入，计算下一状态。

这个模型非常重要，因为 ISA 也可以被理解为状态机规则。以 \instruction{add rax, rbx} 为例，抽象语义是：

\begin{center}
\begin{minipage}{0.72\linewidth}
\begin{verbatim}
RAX_next = RAX_old + RBX_old
RFLAGS_next = flags produced by the addition
\end{verbatim}
\end{minipage}
\end{center}

真实 CPU 内部可能把这条指令拆成微操作，可能乱序执行，可能重命名寄存器，但最终对软件可见的结果必须像上面的状态转移一样。这就是后面“架构状态”和“微架构实现”的区别。

\section{ISA：软件和硬件之间的合同}

\term{ISA}{Instruction Set Architecture} 是指令集体系结构。它定义软件能看到的机器模型：有哪些寄存器、有哪些指令、指令如何编码、内存寻址如何表达、异常如何发生、哪些状态对程序可见。

x86-64 ISA 给程序员承诺了很多东西，例如：

\begin{itemize}
  \item 有一组通用寄存器，如 \register{rax}、\register{rbx}、\register{rcx}、\register{rdx}、\register{rdi}、\register{rsi}、\register{rsp}、\register{rbp}、\register{r8} 到 \register{r15}。
  \item 指令可以读取寄存器、写寄存器、读取内存、写内存、改变控制流。
  \item 内存可以按字节寻址。
  \item 常见寻址模式可以写成 \code{base + index * scale + displacement}。
  \item \register{rip} 表示当前执行位置，\register{rsp} 表示栈顶。
\end{itemize}

但 ISA 不告诉你 CPU 内部有多少流水线阶段，也不告诉你某条指令占用哪个执行端口，不告诉你 L1 cache 多大，不告诉你分支预测器怎么猜。这些属于微架构。

\begin{warningbox}
初学者经常把 ISA 和 CPU 实现混在一起。说“x86-64 有 \register{rax}”是 ISA 层；说“某 Intel 核心上整数加法吞吐约为每周期若干条”是微架构层。优化时必须区分这两层，否则会把一台机器上的性能经验误当成所有 x86-64 机器的规律。
\end{warningbox}

\section{微架构：同一份 ISA 可以有不同的执行机器}

同样执行 x86-64 指令，不同 CPU 的内部结构可以非常不同。一个简单核心可以顺序执行；一个高性能核心会取指、解码、分派、重命名、乱序执行、访存、提交。CPU 还会预测分支、预取数据、维护多级缓存、处理 TLB、解决内存一致性问题。

微架构的存在解释了一个看似奇怪的事实：两段汇编指令数量相同，性能可能差很多；两段 C++ 源码看起来复杂度相同，真实性能也可能差很多。

例如：

\begin{lstlisting}
add eax, dword ptr [rdi + rcx*4]
\end{lstlisting}

这条指令表面上是“从数组加载一个 int 并加到 eax”。微架构层至少可能发生：

\begin{itemize}
  \item 地址生成单元计算 \code{rdi + rcx * 4}。
  \item load 单元向 L1 data cache 发起读取。
  \item 如果虚拟地址还没有 TLB 命中，需要页表相关机制参与。
  \item 如果 L1 miss，访问 L2；L2 miss，访问 L3；再 miss，访问内存。
  \item 加法依赖 load 的结果，因此 load latency 会影响后续加法。
  \item 如果循环中有多个独立 load，乱序执行可以重叠等待。
\end{itemize}

同一条汇编，在 cache 命中和 cache miss 时的代价完全不同。这就是为什么本书必须把“指令”和“数据所在位置”一起讲。

\section{操作系统：程序不是独占机器}

C++ 程序不是直接裸奔在 CPU 上。普通 Linux 程序运行在操作系统提供的进程环境中。操作系统负责创建进程、建立虚拟地址空间、加载可执行文件、处理系统调用、调度线程、管理页面、处理中断和异常。

从性能角度看，操作系统至少带来五类重要影响：

\begin{enumerate}
  \item 虚拟内存让程序看到的是虚拟地址，不是物理地址。
  \item 页面和 TLB 会影响大数组、随机访问和多线程程序。
  \item 调度会让线程在不同核心之间迁移，影响 cache 热度和测量稳定性。
  \item 系统调用和 page fault 会造成远高于普通指令的开销。
  \item CPU 频率、电源管理、后台任务和中断会污染 benchmark。
\end{enumerate}

因此，真正的性能工程不是“写一段快代码”这么简单。你还要控制运行环境，知道什么时候该固定 CPU 频率，什么时候该 pin 线程，什么时候该丢弃第一次运行，什么时候要看 page fault、context switch 和 PMU 事件。

\section{编译器：它不是翻译员，而是证明器和重写器}

很多初学者把编译器理解成“把 C++ 翻译成汇编的工具”。这个理解只对了一半。现代优化编译器更像一个在语言规则约束下工作的证明器和重写器。

它会问：

\begin{itemize}
  \item 这个变量的值是否一定不被使用？
  \item 这个循环是否可以展开？
  \item 两个指针是否可能指向同一块对象？
  \item 这个函数是否可以内联？
  \item 这个条件是否永远成立或永远不成立？
  \item 这个浮点表达式能不能重排？
  \item 这个内存访问是否连续，能不能向量化？
\end{itemize}

只要 C++ 标准允许，编译器就可能把源码改写成完全不同的机器执行形式。比如下面的代码：

\begin{lstlisting}[language=C++]
int f(int x)
{
    return x + 1 > x;
}
\end{lstlisting}

在没有有符号整数溢出的前提下，编译器可以认为 \code{x + 1 > x} 恒为真，从而返回 1。可是如果你按机器补码回绕去想，\code{x} 等于最大正整数时似乎会溢出。关键在于：C++ 有符号整数溢出是未定义行为，编译器可以假设它不会发生。

\begin{deepdive}
性能优化必须尊重语言规则。很多“我看机器码应该这样”的推理，如果违反 C++ 对对象生命周期、别名、对齐、未定义行为的规定，就不是可靠优化，而是在给编译器提供错误前提。后续讲 alias、\code{restrict}、\code{std::bit_cast}、fast-math 和自动向量化时，会反复回到这一点。
\end{deepdive}

\section{ABI：函数之间、模块之间必须说同一种话}

\term{ABI}{Application Binary Interface} 是二进制接口。它规定函数调用时参数如何传递、返回值放在哪里、哪些寄存器由调用者保存、哪些由被调用者保存、栈如何对齐、结构体如何返回、异常和动态链接如何约定。

源码层面两个函数看起来只是：

\begin{lstlisting}[language=C++]
int g(int x, int y);
int z = g(1, 2);
\end{lstlisting}

机器层面却必须回答：

\begin{itemize}
  \item \code{1} 放进哪个寄存器？
  \item \code{2} 放进哪个寄存器？
  \item 调用前 \register{rsp} 要满足什么对齐要求？
  \item \code{g} 能随便改哪些寄存器？
  \item 返回值从哪里取？
\end{itemize}

System V AMD64 ABI 是 Linux x86-64 上的核心规则之一。后面学习汇编时，如果不懂 ABI，就看不懂函数序言、函数尾声、栈帧、寄存器保存和 C++ 与汇编互操作。

\section{内存：性能工程的第二个主角}

第一个主角是指令，第二个主角是数据。现代 CPU 的算术单元非常快，但内存层级相对慢得多。一个 AI 算子是否快，常常不取决于“乘法会不会写”，而取决于数据如何布局、是否连续、是否复用、是否落在 cache、是否适合 SIMD load/store。

看下面两个结构：

\begin{lstlisting}[language=C++]
struct ParticleAoS {
    float x;
    float y;
    float z;
    float mass;
};

struct ParticleSoA {
    float* x;
    float* y;
    float* z;
    float* mass;
};
\end{lstlisting}

AoS 和 SoA 的差异不是语法审美问题，而是内存访问模式问题。如果一个 kernel 只需要所有粒子的 \code{x}，AoS 会把不需要的 \code{y}、\code{z}、\code{mass} 一起拖进 cache line；SoA 则能连续加载 \code{x}。后续讲 SIMD、cache line、GEMM packing 和 AI 算子 fusion 时，这个思想会不断出现。

\section{性能不是运行一次得到的数字}

性能工程最危险的习惯之一，是把一次运行时间当成真相。真实机器上，一次运行结果可能受到很多因素影响：

\begin{itemize}
  \item 编译器是否把代码优化掉。
  \item 输入数据是否还在 cache 中。
  \item CPU 当前频率是多少。
  \item 线程是否被调度到另一个核心。
  \item 后台进程和中断是否打断了运行。
  \item 分支预测器和 TLB 是否已经被预热。
  \item 第一次执行是否触发动态链接、page fault 或内存分配。
\end{itemize}

所以本书要求你形成“证据链”意识。一个高质量性能结论通常至少包含：

\begin{enumerate}
  \item correctness reference：优化前后结果是否一致。
  \item 编译参数：是否启用合理优化，是否说明目标 ISA。
  \item 汇编或 IR：机器实际执行的形态是什么。
  \item benchmark 方法：如何预热、重复、统计、避免被优化删除。
  \item profiling 或 PMU：瓶颈证据支持哪个解释。
  \item 反例检查：有没有输入规模或机器配置让结论失效。
\end{enumerate}

\begin{warningbox}
“我的版本快了 20\%”不是结论，只是观察。真正的结论必须说明：在哪台机器、什么编译器、什么参数、什么输入规模、什么统计方法下快了 20\%；为什么快；有什么证据排除了测量噪声和错误优化；在哪些条件下不保证快。
\end{warningbox}

\section{正确性是性能的前提}

性能优化有一个简单但经常被忽略的底线：错误程序没有性能意义。一个函数快，但结果错了；一个 benchmark 快，但热路径被编译器删除了；一个手写汇编快，但破坏了调用约定，这些都不能算优化成功。

正确性在本书里至少有四层。

第一层是语言正确性。C++ 程序不能依赖未定义行为。越界访问、use-after-free、未初始化读取、错误对齐、违反严格别名、数据竞争、有符号整数溢出等问题，都可能让编译器做出看似奇怪但合法的优化。性能代码越接近底层，越容易触碰这些边界。

第二层是功能正确性。优化前后结果要一致。对整数函数，通常可以逐值比较；对浮点函数，需要定义误差标准；对随机算法，需要固定种子或比较统计性质；对近似算法，需要说明允许误差来自哪里。

第三层是二进制接口正确性。函数调用必须遵守 ABI。参数寄存器、返回值、栈对齐、callee-saved 寄存器、异常边界和符号名都不能随意处理。手写汇编或跨语言互操作时，功能测试通过还不够，因为某些 ABI 破坏可能只在特定调用者、优化级别或异常路径下暴露。

第四层是测量正确性。计时区域必须包含想测的工作，不包含不该测的工作；结果必须被使用，防止死代码消除；输入必须代表目标场景；重复次数和统计口径必须清楚。错误测量会把优化方向带偏。

\begin{center}
\begin{tabular}{p{0.22\linewidth}p{0.38\linewidth}p{0.26\linewidth}}
\toprule
正确性层次 & 常见失败 & 主要证据 \\
\midrule
语言正确性 & 未定义行为、对象生命周期错误 & 警告、sanitizer、代码审查 \\
功能正确性 & 输出错误、误差超界 & reference、单元测试、随机测试 \\
ABI 正确性 & 寄存器破坏、栈不对齐、符号不匹配 & ABI 推导、GDB、反汇编 \\
测量正确性 & 空循环、测错区域、样本污染 & harness 审计、原始数据、报告 \\
\bottomrule
\end{tabular}
\end{center}

这四层并不是最后才检查。正确流程是先有 reference，再写 baseline，再做优化；每次修改都重新验证正确性和二进制形态。很多高水平工程师并不是因为知道更多指令才稳定，而是因为他们不会让错误结果混进性能比较。

\section{从观察到结论需要中间推理}

本书会反复要求你区分观察、解释和结论。

观察是工具或实验直接给出的东西。比如反汇编里出现 \instruction{cmov}，GDB 显示 \register{rsp} 某个值，CSV 中某组样本中位数更低。

解释是你用模型把观察连起来。比如因为 \instruction{cmov} 消除了数据相关分支，所以在随机输入下可能减少 branch miss；因为 \register{rsp} 没有按 ABI 对齐，所以某些调用可能崩溃或变慢；因为样本分布有长尾，所以均值比中位数更容易受噪声影响。

结论是解释在证据约束下能站住的说法。它应该有边界。比如“在这个输入分布和这台机器上，条件移动版本比分支版本更稳定；尚未证明在有强偏置分支时仍然更快”。这比“cmov 一定比 branch 快”更专业。

一个完整推理链可以写成：

\begin{enumerate}
  \item 源码层：两个版本功能等价，并通过 reference 校验。
  \item 编译层：二者使用相同编译器和优化级别。
  \item 二进制层：版本 A 使用条件跳转，版本 B 使用条件移动。
  \item 输入层：输入分布接近随机，分支方向难预测。
  \item 测量层：多次样本显示版本 B 中位数更低，方差更小。
  \item 解释层：较少分支错误可能是主要原因。
  \item 限制层：还需要 \code{perf stat} 的 branch-misses 支持，并测试偏置输入。
\end{enumerate}

清楚的解释不是堆更多词，而是把每一步证据和假设放在正确位置。

\section{从代码到证据的基本流程}

底层分析不能从“这个写法更快”“这个指令更高级”开始，而要从可复查对象开始。一个可复查的优化流程至少包含下面十步：

\begin{enumerate}
  \item 写出清晰正确的 baseline。
  \item 准备 correctness reference 和误差标准。
  \item 建立粗略性能模型：计算量、访存量、预期瓶颈。
  \item 设计 benchmark，避免死代码消除和输入偏差。
  \item 编译并查看汇编或 IR。
  \item 运行测量，记录环境和统计结果。
  \item 用 profiling 或 PMU 验证瓶颈。
  \item 修改一个因素：布局、循环、分块、SIMD、预取或并行。
  \item 再次验证正确性和性能。
  \item 写报告，说明哪些假设被证实，哪些被推翻。
\end{enumerate}

这套流程的核心是顺序。没有 reference，速度数字没有意义；没有二进制证据，就不知道机器实际执行了什么；没有原始样本，就无法判断一次结果是否只是噪声；没有限制条件，结论就会被滥用。

同样，对一段 C/C++ 代码做底层阅读时，也应按固定顺序拆开：输入和副作用是什么，依赖哪些 C++ 语义前提，热路径在哪里，循环归纳变量和退出条件是什么，每次迭代访问哪些地址，关键值落在寄存器、栈还是内存里，是否存在循环携带依赖，分支是否与数据分布相关，编译器能做哪些合法重写，最后需要哪些工具证据支撑判断。

以 \code{sum_i32} 为例，输入是数组地址和长度，输出是累加和；每次迭代连续读取一个 32 位整数；累加变量形成依赖；循环控制产生分支；编译器可能展开或向量化；如果结果没有被使用，整个计算可能被删除。要把这些说法变成结论，必须看优化后的汇编和测量样本，而不是只看源码。

\section{七个合同：理解系统分层的核心工具}

前面说过计算机系统要分层。更精确地说，层与层之间依靠合同连接。合同规定了上层可以依赖什么，下层可以自由实现什么。学习底层时，最重要的能力之一就是知道自己正在使用哪个合同。

\subsection{语言合同}

C++ 源码首先受语言合同约束。语言合同规定表达式如何求值、对象什么时候存在、指针能指向哪里、引用是否必须绑定有效对象、整数溢出有什么规则、浮点运算是否允许重排、线程之间怎样形成数据竞争。编译器优化的出发点就是语言合同，而不是某台机器上的直觉。

例如越界访问在某次运行中可能“碰巧读到一个值”，但语言合同不保证这种行为。编译器可以利用“程序不会越界”这个前提重排或删除代码。再例如一个对象生命周期结束后，旧指针的数值可能还像地址，但语言层已经不允许通过它访问对象。很多性能 bug 的根源不是硬件太复杂，而是程序先违反了语言合同，然后编译器在合法范围内做了读者没有预期的优化。

\subsection{编译合同}

编译器接收源程序，输出等价的目标程序。这里的等价不是“每行源码对应几条汇编”，而是“对语言定义的可观察行为等价”。只要可观察行为不变，编译器可以内联、删除、合并、常量传播、循环展开、向量化、重排，甚至让你在汇编里找不到原来的局部变量。

这就是为什么调试构建和发布构建表现差异巨大。调试构建通常保留更多局部变量和栈帧形态，方便观察；优化构建更关心机器执行。读者必须学会在两种构建之间切换：用调试构建理解源码意图，用优化构建分析真实性能。

\subsection{目标文件和链接合同}

一个翻译单元编译后通常先变成目标文件。目标文件包含机器码、符号、重定位信息、节区和调试信息。链接器再把多个目标文件、静态库和动态库组合成可执行文件或共享库。这里的合同决定函数名如何被解析、外部符号在哪里、全局对象如何初始化、代码和数据放进哪些段。

如果只看单个源文件，就会错过很多二进制事实。函数可能在另一个目标文件里；库函数可能通过动态链接进入；全局符号可能因为可见性和重定位产生额外间接层；链接时优化可能跨文件改写调用关系。后面第 1 章和第 4 章会把这个过程具体化。

\subsection{ABI 合同}

ABI 合同让不同编译单元、不同语言、不同汇编文件可以互相调用。它规定参数寄存器、返回值寄存器、栈对齐、寄存器保存责任、结构体传参、异常边界和符号命名。ABI 的重要性在手写汇编时尤其明显：只要你破坏 callee-saved 寄存器或栈对齐，程序可能在当前测试里没崩，但在另一个调用点突然出错。

因此，汇编分析的重点不是背几条指令，而是看清二进制接口的责任边界：调用者保证什么，被调用者保证什么，返回后哪些状态仍然有效，哪些状态可以被破坏。

\subsection{操作系统合同}

普通程序运行在操作系统提供的进程合同里。进程拥有虚拟地址空间、文件描述符表、当前工作目录、环境变量、权限、信号处理方式和线程集合。操作系统负责把可执行文件加载进内存，给用户态程序提供系统调用入口，在异常和中断发生时接管控制。

这个合同解释了为什么同一个二进制在不同环境下表现不同。当前目录改变，打开文件可能失败；环境变量改变，动态库搜索路径可能改变；地址空间随机化改变，某些地址数值会变；内存压力改变，page fault 可能增多；调度策略改变，benchmark 波动可能变大。

\subsection{ISA 合同}

ISA 合同是软件可见机器状态的规则。它规定 \register{rax}、\register{rsp}、\register{rip} 和标志位这样的架构状态如何被指令读取和修改。汇编阅读的核心就是在 ISA 合同下做状态追踪。看到 \instruction{cmp}，要知道它主要写标志位；看到 \instruction{jle}，要知道它读取标志位决定控制流；看到 \instruction{call}，要知道它写返回地址并改变 \register{rip}。

ISA 合同不承诺内部执行成本。它告诉你指令“应该产生什么结果”，不告诉你“用了多少周期”。后者属于微架构和具体机器状态。

\subsection{测量合同}

测量合同是你自己给实验制定的规则。它规定什么算输入、什么算输出、什么代码在计时区域内、重复多少次、如何统计、如何保存原始数据、如何证明 correctness。很多性能争论没有结果，就是因为双方没有共享同一个测量合同：一个人在测端到端，一个人在测热循环；一个人在测 cold cache，一个人在测 warm cache；一个人报告最小值，一个人报告均值。

合同意识必须从第一章开始建立。读任何性能文章、benchmark 图表或优化报告，都先问它的合同是什么；没有说清的部分，就是结论最脆弱的部分。

\section{必要背景：从语义到表示}

进入底层分析前，需要把 C++ 语义、二进制表示、工具和基础数学接起来，而不是把它们当作互不相关的预备知识。变量有类型、对象身份、生命周期、存储期和可能的地址；指针不只是一个整数地址，还受对象边界、对齐、别名和生命周期限制；引用不是可空指针，数组到指针转换不是数组复制，函数调用也不是文本替换。未定义行为不是“随机返回一个值”，而是语言合同已经不再给出保证。

这些规则直接影响汇编形态。一个越界访问在某次运行中可能读到数值，但编译器仍可假设合法程序不会越界；一个局部变量在 \code{-O3} 下可能没有固定栈槽，因为它被寄存器化、常量传播或完全删除；一个 \code{const} 指针参数也不等于底层物理内存只读，它只是限制通过该表达式进行修改。可靠的性能解释必须同时尊重语言语义和机器成本。

\subsection{二进制和整数}

机器用位表示数据。一个 32 位模式可以解释成无符号整数、有符号整数、浮点数、指令编码、地址的一部分，也可以只是字节序列。表示和解释必须分开。

例如位模式 \code{0xffffffff} 可以表示无符号整数最大值，也可以在常见补码机器上表示有符号 \code{-1}。但是 C++ 的有符号溢出规则、无符号回绕规则和机器补码表示不是同一个层次。后续学习汇编时，你会看到同一条 \instruction{add} 可以用于有符号和无符号加法；区别常常体现在你如何解释标志位和后续跳转，而不是加法器本身换了一套。

需要掌握的基础包括：

\begin{itemize}
  \item 位、字节、字、十六进制表示。
  \item 补码的基本思想。
  \item 左移、右移、掩码、按位与或异或。
  \item 大端和小端的概念。
  \item 整数宽度和截断。
  \item 地址计算中的 scale 和位移。
\end{itemize}

这些知识看似低级，但会直接影响你是否看得懂反汇编。看到 \code{[rdi + rsi*4 + 8]}，你要立刻想到“基址、索引、元素大小、字段偏移”，而不是把它当成一串难读符号。

\subsection{Linux 工具链}

本书默认在 Linux 或接近 Linux 的环境中学习。原因不是其他系统不重要，而是 Linux 工具链更容易直接观察编译、链接、进程、内存映射、系统调用和性能计数。你不需要成为运维专家，但必须能在 shell 中完成基本实验。

至少要熟悉：

\begin{itemize}
  \item 文件路径、当前目录、相对路径和绝对路径。
  \item \code{mkdir}、\code{cp}、\code{mv}、\code{rm}、\code{find}、\code{grep} 或 \code{rg}。
  \item 重定向、管道、退出码和环境变量。
  \item \code{gcc}、\code{clang}、\code{make} 的基本使用。
  \item \code{objdump}、\code{readelf}、\code{nm}、\code{ldd} 的基本用途。
  \item \code{gdb}、\code{strace}、\code{perf} 的基本概念。
\end{itemize}

工具不是附属品。底层学习的很多知识必须通过工具验证。你不能只说“编译器可能内联”，你要用反汇编证明；不能只说“程序访问了非法地址”，你要用 GDB、sanitizer 或 core dump 证明；不能只说“测量有噪声”，你要保存样本并展示分布。

\subsection{数学和算法}

本书不要求读者提前学完高等数学，但一些基础必须准确。算法复杂度要能用于规模分析；对数和指数要能理解数量级；均值、中位数、分位数、方差要能用于读 benchmark；简单线性模型要能用于估算吞吐和带宽。

更重要的是把数学语言和机器现象连接起来。复杂度里的 $O(n)$ 只说明增长阶，机器成本还要乘上每个元素的字节数、每次迭代的指令数、依赖长度、cache miss 概率和并行度。均值是样本摘要，不是全部事实；中位数能代表典型情况，但可能隐藏长尾；最小值常被用来近似低噪声状态，但不能代表用户体验。

学习性能时，数学不是为了形式化炫技，而是为了防止语言含糊。说“快很多”不如说“中位数降低 18\%，最小值降低 15\%，p95 没有改善”。说“内存访问很多”不如说“每个元素读取 8 字节、写 8 字节，理论最少 16n 字节流量”。这种表达会迫使你面对事实。

底层分析中最坏的习惯是过早确定答案。程序崩溃时，不要先说“肯定是栈坏了”；benchmark 波动时，不要先说“肯定是系统调度”；汇编变短时，也不能直接说“肯定更快”。正确做法是列出假设、设计观察、逐步排除。

例如一个函数变慢，可以先列出：

\begin{enumerate}
  \item 编译器生成了不同机器码。
  \item 输入规模改变导致 cache 状态不同。
  \item 分支预测命中率下降。
  \item 内存分配或初始化混进计时区域。
  \item CPU 频率或调度状态改变。
  \item correctness 校验路径不一致。
\end{enumerate}

然后逐项找证据。看反汇编验证第一个假设；固定输入和规模验证第二个；用 \code{perf stat} 或输入分布验证第三个；审计 harness 验证第四个；记录环境和 raw samples 验证第五个；统一 reference 验证第六个。这个过程比经验猜测慢一点，但它能排除错误解释。

\section{把一个表达式拆到机器层}

为了把抽象链路落地，来看一个更小的表达式：

\begin{lstlisting}[language=C++]
std::uint64_t y = a[i] + 7;
\end{lstlisting}

源码层说：从数组中取第 \code{i} 个元素，把它和常量 7 相加，得到 \code{y}。但底层分析要继续拆。

语言层首先问：\code{a} 是否指向至少 \code{i + 1} 个有效元素？\code{i} 是否在范围内？\code{a[i]} 的类型是什么？读取是否违反对象生命周期？如果这些不成立，程序语义已经不可靠。

类型和布局层问：元素宽度是多少？如果 \code{a} 指向 \code{std::uint64_t}，每个元素 8 字节；如果指向 \code{std::uint32_t}，每个元素 4 字节。数组下标不是神秘语法，本质上会引出地址计算。

编译层问：\code{i} 是否已知？\code{a} 是否可能和 \code{y} 所在对象别名？表达式是否在循环中？编译器能否把地址计算和 load 移动、合并或删除？如果 \code{y} 后续没有可观察使用，整句可能被删掉。

ISA 层问：基址在哪个寄存器？索引在哪个寄存器？使用哪种寻址模式？加法写哪个寄存器？标志位是否重要？可能的形态是：

\begin{lstlisting}
mov rax, qword ptr [rdi + rsi*8]
add rax, 7
\end{lstlisting}

微架构层问：这个 load 是否命中 L1？地址生成是否依赖前面的计算？后面的指令是否依赖 \register{rax}？如果在循环中，是否有多个独立 load 可以重叠？如果每次 \code{i} 来自另一个数组，硬件预取是否还能工作？

测量层问：如果你想比较两种写法，计时区域是否只包含表达式所在 kernel？输入是否防止编译器常量折叠？结果是否被消费？样本是否重复？是否看过汇编确认两种写法真的生成不同代码？

\begin{mentalmodel}
底层学习的核心动作是“拆开”。把一行源码拆成语义前提、类型布局、地址计算、指令状态、微架构成本和测量证据。只有拆到这些层次，优化讨论才不会停留在表面。
\end{mentalmodel}

\section{十二章之间的依赖关系}

第一册虽然只有 12 章，但不是 12 篇孤立文章。它们之间有清楚的依赖。

第 1 章解释程序如何从源码进入进程世界。没有这章，就容易把源码、目标文件、可执行文件和运行中进程混成一件事。第 2 章解释 CPU 如何执行指令，帮助读者建立架构状态、依赖、分支和访存的基本模型。第 3 章解释数据、内存、指针和布局，把 C++ 对象和机器地址连接起来。第 4 章讲 Linux 工具和调试方法，让前面三章的概念能被真实观察。

第 10 章开始进入 x86-64 汇编语法。它不是突然跳到另一门语言，而是把前面“机器状态”具体化。第 11 章讲整数指令，让读者能追踪算术、位运算和标志位。第 12 章讲控制流，让读者理解分支、循环、比较和函数跳转。第 13 章讲 System V ABI，把函数调用变成可验证的二进制合同。第 14 章讲 C++ 与汇编互操作，要求读者真正把 ABI、类型布局、栈对齐和测试结合起来。

第 15 章和第 16 章把前面的观察能力转成测量能力。第 15 章问“性能数字为什么可能不可信”，第 16 章问“怎样设计一个长期可维护的 benchmark”。到这里，第一册形成闭环：你能读源码，能看二进制，能理解执行状态，能证明接口正确，能设计测量，并能写出有边界的结论。

\begin{center}
\begin{tabular}{p{0.18\linewidth}p{0.34\linewidth}p{0.34\linewidth}}
\toprule
阶段 & 核心问题 & 产出能力 \\
\midrule
第 0 到 4 章 & 程序如何成为进程和机器状态。 & 能观察源码、进程、内存和指令。 \\
第 10 到 12 章 & 汇编如何表达计算和控制流。 & 能追踪寄存器、标志位、跳转和循环。 \\
第 13 到 14 章 & 函数边界如何在二进制层成立。 & 能写和审计 C++ 与汇编互操作。 \\
第 15 到 16 章 & 性能数字如何成为可信证据。 & 能设计、运行和解释 benchmark。 \\
\bottomrule
\end{tabular}
\end{center}

学习时不要跳过前面的观察训练。很多人急着看汇编，却不知道可执行文件如何生成；急着看 benchmark，却不知道优化器实际删除了什么；急着写手工汇编，却不知道 ABI 保存责任。第一册的章节顺序就是为了避免这些断裂。

\section{从第一册走向第二册}

第一册结束时，读者应当具备进入第二册的条件。第二册可以把主线扩展到更深的性能主题：cache 层级、带宽模型、SIMD、循环优化、GEMM、AI 算子、PMU、线程、NUMA 和生产级性能诊断。那时每个高级主题都可以接到第一册的基础链路上。

例如讲 cache line 时，可以回到第 3 章的对象布局和地址；讲 SIMD 时，可以回到第 10 到 12 章的寄存器和指令；讲手写 microkernel 时，可以回到第 13 到 14 章的 ABI；讲 PMU 时，可以回到第 15 到 16 章的测量合同；讲 AI 算子时，可以回到算法复杂度、访存字节数、数值正确性和 benchmark 报告。

\section{把底层知识连成因果链}

学习底层系统时，最重要的不是记住更多术语，而是能把术语连成因果链。一个事实如果不能进入因果链，就很难用于解释真实问题。例如“CPU 有 cache”只是孤立知识；“数组连续访问通常比随机访问更容易利用 cache line，因此在同样算法复杂度下可能有更高吞吐”才是可以用于分析的因果链。再进一步，“这个 benchmark 的工作集从 32 KiB 增长到 64 MiB 后吞吐下降，可能是因为工作集越过 cache 容量；需要用输入规模扫描和 PMU 事件验证”就是工程化因果链。

本书要求你在每个概念后面追问五个问题：

\begin{enumerate}
  \item 它属于哪一层：语言、编译器、ABI、ISA、操作系统、微架构还是测量。
  \item 它约束了谁：程序员、编译器、链接器、加载器、CPU 还是测量者。
  \item 它能解释什么现象：崩溃、错误结果、汇编形态、性能差异还是工具输出。
  \item 它需要什么证据：源码、命令、反汇编、调试器、计数器、样本数据还是规范文档。
  \item 它有什么边界：在哪些平台、优化级别、输入范围或实验协议下成立。
\end{enumerate}

这五个问题能把“知道一个词”变成“会使用一个概念”。比如“ABI”这个词，如果只解释成“调用约定”，还不够。你要知道它属于函数和模块边界，约束调用者和被调用者，能解释参数寄存器、返回值、栈对齐、callee-saved 保存、名字修饰和动态库兼容，需要通过函数声明、反汇编、符号表和 ABI 文档验证，边界则包括平台、编译器、语言类型和链接模式。

\subsection{因果链要允许被推翻}

工程分析中的因果链不是作文，而是可被证据推翻的假设。你可以先假设“这个循环慢是因为随机访问导致 cache miss”，但随后要设计证据：改变访问模式、扫描工作集大小、查看反汇编、收集 cache miss 相关事件、比较顺序访问 baseline。若证据不支持，就要修改假设。不能因为某个解释听起来高级，就把它固定成结论。

高质量学习笔记应保留这种修正过程。例如：

\begin{lstlisting}
Initial hypothesis:
    random access is slower because of cache misses

Evidence:
    sequential and random versions have similar objdump shape
    random version slows down only when data exceeds L2-sized range
    perf stat shows more cache misses for large sizes

Revised conclusion:
    cache behavior explains the large-size gap, but small-size gap is likely
    dominated by loop overhead and branch/data dependency differences
\end{lstlisting}

这样的笔记比直接写“随机访问 cache miss 多”更有价值，因为它显示了证据如何约束结论。

\section{第一册的知识地图}

第一册虽然只有十二章，但覆盖了从源码到可信测量的完整基础链路。可以把它看成六个连续区域。

第一区域是系统总图。第 0 章建立抽象层和学习方法，第 1 章把程序从源码带到进程，第 2 章解释 CPU 如何把指令作为状态变化执行。这一区域回答“程序如何成为正在运行的东西”。

第二区域是数据和工具。第 3 章解释对象、内存、指针、布局、对齐和别名，第 4 章解释 Linux 工具、构建、调试和报告。这一区域回答“程序状态如何被表示、观察和复现”。

第三区域是汇编阅读。第 10 章讲汇编语法、符号、section、寻址和机器码，第 11 章讲整数指令、flags、扩展、位运算和地址计算，第 12 章讲控制流、循环、调用和 switch。这一区域回答“如何从机器码恢复数据流和控制流”。

第四区域是二进制接口。第 13 章讲 System V ABI，第 14 章讲 C++ 与汇编互操作。这一区域回答“函数、模块、语言边界如何共享同一个二进制合同”。

第五区域是测量。第 15 章讲可信测量的基本原则，第 16 章讲 benchmark suite 设计。这一区域回答“如何把运行时间变成可审查证据，而不是一次性数字”。

第六区域是后续第二册的入口。AI/HPC/SIMD/微架构优化不在第一册展开，但第一册为它们准备必要前提：能读汇编，能解释 ABI，能分析内存布局，能设计 benchmark，能区分事实、推测和边界。

\section{综合案例：一次性能问题应该怎样被拆开}

假设你接到一个问题：“新版本比旧版本慢了 20\%。”这句话本身不是一个可处理的问题，因为它没有说明哪一个程序、哪一个输入、哪一个构建、哪一个指标、哪一台机器、哪一段路径。第一册训练的正是把这种模糊说法拆成可验证问题。

第一步，确认对象。新版本和旧版本分别对应哪个 Git 提交？工作区是否有未提交修改？运行的是哪个可执行文件？动态库是否来自同一个目录？编译器和参数是否一致？如果这些问题没有回答，所谓“新旧版本”只是口头标签，不是工程对象。

第二步，确认正确性。两个版本是否对同一输入产生同一结果？是否使用同一个 reference？是否覆盖边界输入？如果新版本少处理了尾部、跳过了错误检查、使用了未定义行为，它可能更快或更慢，但这种性能比较没有意义。正确性是性能的前提，不是性能之后才补的手续。

第三步，确认测量协议。计时区域包含什么？是否把输入生成、内存分配、打印、校验放进 timed region？重复次数和 warmup 是多少？raw samples 是否保存？如果只有一次运行时间，就只能说观察到一次结果，不能说形成性能结论。

第四步，确认二进制。热点函数是否还存在？是否被内联、展开、向量化、删除、替换成库调用？当前机器是否走了某个 ISA dispatch 路径？如果没有反汇编或符号证据，就不知道测到的是源码意图、优化后热路径，还是完全不同的实现。

第五步，建立假设。若新版本多了随机访问，可能是 cache 或 TLB；若多了间接调用，可能是内联失败或分支目标预测；若多了小对象分配，可能是 allocator 和 page fault；若汇编变成标量，可能是别名信息阻止向量化。每个假设都必须能对应下一步证据。

第六步，设计对照。改变输入规模可以区分固定开销和带宽瓶颈；改变数据分布可以区分分支预测和算术成本；固定 CPU 核心可以减少调度噪声；比较 objdump 可以确认机器码变化；使用 \code{perf stat} 可以为分支、cache 或指令数假设提供硬件证据。

最后，写出有边界的结论。合格结论不是“新版本慢 20\%”，而是类似：“在这台机器、这个编译器、这个输入规模和 warm-cache 协议下，新版本 median ns/element 增加 20\%。正确性检查通过。反汇编显示新版本热循环多了一次间接 load，规模扫描显示大工作集退化更明显。当前假设是内存访问局部性变差；尚未在另一台机器复测。”

\begin{keyidea}
底层分析的基本动作是把一句模糊现象拆成对象、正确性、协议、二进制、假设、对照和结论。这个动作贯穿第一册所有章节。
\end{keyidea}

\section{本章实验：建立你的机器档案}

\begin{labbox}
本实验不是为了追求速度，而是为了建立后续所有实验都要引用的机器档案。请在 \filepath{results/machine-profile/} 下保存报告。

\subsection*{任务 1：记录 CPU 和系统信息}

运行：

\begin{lstlisting}[language=bash]
lscpu
cat /proc/cpuinfo | sed -n '1,80p'
cat /proc/meminfo | sed -n '1,40p'
uname -a
gcc --version
clang --version
ldd --version
\end{lstlisting}

报告中至少记录：

\begin{itemize}
  \item CPU 型号、核心数、线程数。
  \item L1d/L1i/L2/L3 cache 信息。
  \item 支持的 ISA 扩展，例如 SSE、AVX、AVX2、AVX-512、FMA、VNNI。
  \item 内存总量。
  \item Linux kernel 版本。
  \item GCC/Clang 版本。
\end{itemize}

\subsection*{任务 2：建立第一个源码到汇编证据链}

创建 \filepath{labs/ch00_machine_map/sum.cpp}：

\begin{lstlisting}[language=C++]
#include <cstddef>
#include <cstdint>

extern "C" std::int64_t sum_i32(std::int32_t const* a, std::size_t n)
{
    std::int64_t s = 0;
    for (std::size_t i = 0; i < n; ++i) {
        s += a[i];
    }
    return s;
}
\end{lstlisting}

分别生成未优化和优化后的汇编：

\begin{lstlisting}[language=bash]
clang++ -std=c++20 -O0 -S -masm=intel sum.cpp -o sum_O0.s
clang++ -std=c++20 -O3 -S -masm=intel sum.cpp -o sum_O3.s
\end{lstlisting}

报告必须回答：

\begin{enumerate}
  \item \code{-O0} 和 \code{-O3} 的循环形态有什么差异？
  \item 参数 \code{a} 和 \code{n} 分别从哪些寄存器进入函数？
  \item 汇编中哪里体现了 \code{sizeof(std::int32_t) == 4}？
  \item 编译器有没有展开或向量化？证据是什么？
\end{enumerate}

\subsection*{任务 3：第一次测量但不急着下结论}

写一个小 benchmark，测量 \code{n = 1024}、\code{n = 1048576}、\code{n = 67108864} 三种规模。每个规模运行多次，记录最小值、中位数和最大值。报告中必须写明为什么三种规模可能对应不同 cache 行为。
\end{labbox}

\section{本章作业：写出你的第一份性能工程备忘录}

\begin{exercisebox}
提交一份不少于 1500 字的中文报告，题目为“从 C++ 到机器执行的第一张地图”。报告必须包含以下部分：

\begin{enumerate}
  \item 画出从 C++ 源码到 CPU 执行的链路图，并用自己的话解释每一层的职责。
  \item 用 \code{sum_i32} 例子解释源码、汇编、寄存器、内存访问之间的关系。
  \item 说明 ISA 和微架构的区别，并举一个“同一条汇编在不同数据状态下性能不同”的例子。
  \item 说明操作系统为什么会影响 benchmark。
  \item 列出你当前最不熟悉的 10 个概念，并为每个概念写一句你准备如何验证它。
\end{enumerate}

评分重点不是语言华丽，而是证据完整、层次清楚、没有把不同抽象层混在一起。
\end{exercisebox}

\section{验收标准}

本章收束时，应能做到：

\begin{itemize}
  \item 画出源码到机器执行的完整粗粒度链路。
  \item 区分语言语义、编译器优化、ISA、微架构、操作系统和测量方法。
  \item 对一个简单循环提出至少三种可能性能瓶颈。
  \item 用汇编证据回答“编译器实际生成了什么”。
  \item 解释为什么一次运行时间不能直接当成性能结论。
\end{itemize}

后续章节会把这张地图逐层放大：先看源码如何变成进程，再看 CPU 如何执行指令、数据如何落在内存里，随后进入 x86-64 汇编、ABI、C++ 与汇编互操作，最后用可信 benchmark 方法把观察结果写成证据。内存层级、SIMD、GEMM 和 AI 算子属于第二册。
